\section{Introduction}
Suppose two simulations predict the same number of people attending an activity. In one, almost everyone receives the announcement but few are interested. In the other, few people receive it, but most recipients want to participate. A historical attendance count cannot by itself distinguish these explanations. Their difference becomes consequential when a new activity combines forwarding, a waiting list, and overlapping schedules: forwarding changes who knows about the opportunity, waiting-list turnover changes who can obtain a place, and conflicts determine which intentions become attendance. The example applies to resource allocation and collective participation more generally; a campus is an executable instance, not the intended boundary of the research question.

Calibration from aggregate outputs is therefore an inverse problem with potentially non-identifiable mechanisms. Set-based approaches already recognize that multiple parameters can be statistically eligible \citep{bai2022}. The remaining operational question is which additional process observation is useful when collecting every micro-level record is unnecessary or costly. A record can have high behavioral entropy yet reveal little about a future quantity of interest. Conversely, a modest disagreement about viewing an announcement may be valuable if it separates mechanisms that imply different future congestion trajectories. Goal-oriented experimental design provides the appropriate perspective: the value of evidence depends on the prediction it is meant to improve \citep{chaloner1995,attia2018,zhong2026}.

We instantiate this perspective with a finite mechanism bank, explicit environmental settlement, and cached training-rule combination probes. A shared parameter vector controls attention, participation, and waiting-list dropout. Separate process and trajectory checks form a compatible set. A trajectory-based weighting rule then defines a calibrated mixture, and a weighted covariance ranks binary process records by their expected reduction of probe disagreement per unit cost. The same cache supports posterior updates and a distribution over complete paths, so no model retraining or simulator rerun is required inside a query step.

The experimental outcome is more qualified than this construction might suggest. Development selection strongly favors concentrated trajectory weights. The best predictor consequently retains many candidates in its hard set but almost no effective mixture diversity. Its trajectory score improves while its individual-process Brier score worsens. Acquisition ablations have little visible effect because there is little remaining disagreement to resolve. This is not a peripheral failure: it exposes a tension between optimizing aggregate predictive accuracy and preserving useful uncertainty for subsequent evidence acquisition.

Our contributions are threefold. First, we give an explicit, reproducible formulation of process-constrained set calibration and budgeted evidence selection for a multi-agent simulator that separates behavioral intentions from realized outcomes. Second, we implement the binary predictive-variance identity as a batched covariance operation, integrate hypothetical labels for small batches, and state its computational and statistical assumptions. We do not claim the identity or goal-oriented design principle as new. Third, we provide an audited synthetic study separating calibration gains from acquisition gains, including matched-prior controls, activity-level uncertainty, and unknown-rule stress tests. The resulting evidence identifies posterior concentration as a practical failure mode and limits claims that could otherwise be drawn from favorable aggregate scores.

The work should be read as a methodological and diagnostic study. It provides neither a demonstrated state-of-the-art result on SocioVerse nor a result on a MIT real-behavior benchmark. An interface with the required field names does not constitute an evaluated dataset. This distinction also determines the paper's scope: the general formulation is portable, whereas the empirical conclusions concern the constructed system studied here.
